% !TeX encoding = UTF-8
% !TeX program = xelatex
% !TeX spellcheck = en_US

\documentclass[bachelor]{ustcthesis}
% doctor|master|bachelor [academic|professional] [chinese|english] [print|pdf]
% [super|numebers|authoryear]

\title{基于深度学习的金融风控研究 }
\author{郑非}
\major{计算机科学与技术}
\supervisor{李金龙\ 副教授}
\cosupervisor{郑小林\ 副教授}
% \date{二〇一七年五月一日} % 注释掉则为今日
% \professionaltype{专业学位类型}
% \secretlevel{秘密}        % 绝密|机密|秘密，注释本行则不保密
% \secretyear{20}           % 保密年限

\entitle{Deep Learning in Financial Risk Management}
\enauthor{Fei Zheng}
\enmajor{Computer Science and Technology}
\ensupervisor{Asso. Prof. Jinlong Li}
\encosupervisor{Asso. Prof. Xiaolin Zheng}
% \endate{May 1, 2017}      % Today if commented
% \enprofessionaltype{Professional degree type}
% \ensecretlevel{Secret}    % Top secret|Highly secret|Secret


% 加载宏包和配置
\usepackage{graphicx}
\graphicspath{{figures/}}
\usepackage{booktabs}
\usepackage{longtable}
\usepackage[ruled,linesnumbered]{algorithm2e}
\usepackage{siunitx}
\usepackage{amsthm}
\usepackage{hyperref}

\usepackage{xcolor}
\usepackage{listings}
\renewcommand{\lstlistingname}{代码}
\lstset{ %  
  backgroundcolor=\color{white},   % choose the background color; you must add \usepackage{color} or \usepackage{xcolor}  
  basicstyle=\footnotesize,        % the size of the fonts that are used for the code  
  breakatwhitespace=false,         % sets if automatic breaks should only happen at whitespace  
  breaklines=true,                 % sets automatic line breaking  
  captionpos=bl,                    % sets the caption-position to bottom  
  commentstyle=\color{green},    % comment style  
  deletekeywords={...},            % if you want to delete keywords from the given language  
  escapeinside={\%*}{*)},          % if you want to add LaTeX within your code  
  extendedchars=true,              % lets you use non-ASCII characters; for 8-bits encodings only, does not work with UTF-8  
  frame=single,                    % adds a frame around the code  
  keepspaces=true,                 % keeps spaces in text, useful for keeping indentation of code (possibly needs columns=flexible)  
 %  keywordstyle=\color{blue},       % keyword style  
  %language=Python,                 % the language of the code  
  morekeywords={*,...},            % if you want to add more keywords to the set  
  numbers=left,                    % where to put the line-numbers; possible values are (none, left, right)  
  numbersep=5pt,                   % how far the line-numbers are from the code  
  numberstyle=\tiny\color{gray}, % the style that is used for the line-numbers  
  % rulecolor=\color{black},         % if not set, the frame-color may be changed on line-breaks within not-black text (e.g. comments (green here))  
  showspaces=false,                % show spaces everywhere adding particular underscores; it overrides 'showstringspaces'  
  showstringspaces=false,          % underline spaces within strings only  
  showtabs=false,                  % show tabs within strings adding particular underscores  
  stepnumber=1,                    % the step between two line-numbers. If it's 1, each line will be numbered  
  % stringstyle=\color{orange},     % string literal style  
  tabsize=2,                       % sets default tabsize to 2 spaces  
  %title=myPython.py                   % show the filename of files included with \lstinputlisting; also try caption instead of title  
}  

\DeclareRobustCommand\cs[1]{\texttt{\char`\\#1}}
\newcommand\pkg{\textsf}

\renewcommand\vec{\symbf}
\newcommand\mat{\symbf}
\newcommand\ts{\symbfsf}
\newcommand\real{\mathbf{R}}




\begin{document}

% 研究生论文:
%   封面，原创性声明和授权使用声明
%   frontmatter: 摘要，目录，[图、表清单]，[符号说明]
%   mainmatter: 正文章节，参考文献
%   appendix: 附录
%   backmatter: 致谢，已发表论文列表
%
% 本科生论文:
%   封面
%   frontmatter: 致谢，目录，摘要
%   mainmatter: 正文章节，参考文献
%   appendix: 附录

\maketitle
\makestatement

\frontmatter
\input{chapters/my_acknowlegement.tex}

%% !TeX root = ../main.tex

\begin{abstract}
  摘要是论文内容的总结概括，应简要说明论文的研究目的、基本研究内容、 研究方法或
  过程、结果和结论，突出论文的创新之处。摘要中不宜使用公式、图表，不引用文献。
  博士论文中文摘要一般800～1000个汉字，硕士论文中文摘要一般600个汉字。英文摘要的
  篇幅参照中文摘要。

  关键词另起一行并隔写在摘要下方，一般3～8个词，中文关键词间空一字或用分号“；”隔
  开。英文摘要的关键词与中文摘要的关键词应完全一致，中间用逗号“,”或分号“;”隔开。

  \keywords{中国科学技术大学；学位论文；\LaTeX{} 模板；学士；硕士；博士}
\end{abstract}

\begin{enabstract}
  This is a sample document of USTC thesis \LaTeX{} template for bachelor,
  master and doctor. The template is created by zepinglee and seisman, which
  orignate from the template created by ywg. The template meets the
  equirements of USTC theiss writing standards.

  This document will show the usage of basic commands provided by \LaTeX{} and
  some features provided by the template. For more information, please refer to
  the template document ustcthesis.pdf.

  \enkeywords{University of Science and Technology of China (USTC); Thesis;
  \LaTeX{} Template; Bachelor; Master; PhD}
\end{enabstract}

\tableofcontents

\begin{abstract}

\end{abstract}

\begin{enabstract}

\end{enabstract}
% \listoffigures
% \listoftables
%% !TeX root = ../main.tex

\begin{notation}

  \begin{notationlist}{2em}
    \item[$\displaystyle a$] The number of angels per unit area
    \item[$\displaystyle N$] The number of angels per needle point
    \item[$\displaystyle A$] The area of the needle point
    \item[$\displaystyle \sigma$] The total mass of angels per unit area
    \item[$\displaystyle m$] The mass of one angel
    \item[$\displaystyle \sum_{i=1}^n a_i$] The sum of $a_i$
  \end{notationlist}

\end{notation}



% 也可以使用 nomencl 宏包

% \printnomenclature

% \nomenclature{$\displaystyle a$}{The number of angels per unit are}
% \nomenclature{$\displaystyle N$}{The number of angels per needle point}
% \nomenclature{$\displaystyle A$}{The area of the needle point}
% \nomenclature{$\displaystyle \sigma$}{The total mass of angels per unit area}
% \nomenclature{$\displaystyle m$}{The mass of one angel}
% \nomenclature{$\displaystyle \sum_{i=1}^n a_i$}{The sum of $a_i$}


\mainmatter
\chapter{绪论}
\section{课题背景}
一直以来，贷款业务就是商业银行的核心业务之一。但是贷款业务伴随着信用风险。如果贷款者没有按时偿还贷款，银行就会承担经济损失。因此如何尽可能保证贷款者的可靠性，减少坏账的发生率，是银行急需解决的问题之一。同样的需求也存在于提供贷款服务的互联网金融公司。

近年来，随着计算机性能的提高，深度学习在各个领域的应用愈加广泛。在图像识别、机器翻译、模式识别等领域，深度学习的模型效果都达到了业界的顶尖水准\cite{lecun2015deep}。

深度学习可以看做是一个强大的函数拟合工具，特别适用于普通的统计学方法难以处理的超高维数据。对于风险控制的场景来说，数据就是商业银行或互联网金融公司所得到的用户画像，标签就是用户违约的概率。已有的用户画像及其是否违约的记录，就是训练数据。


\section{研究现状}
\subsection{传统的风控方法}
传统的风险控制方法，包括银行工作人员对贷款者的情况进行主观分析，从而判断其偿还能力。银行一般使用的是"The Five Cs of Credit"\cite[122]{apostolik2009foundations},从如下五个维度判断贷款者的信用：
\begin{itemize}
	\item Character: 贷款者的“人格”，包含其名声、能力等。
	\item Capital: 贷款者的资本状况，包含其公司的运营情况、现金流量、杠杆水平、固定资产等。很显然资本状况好的贷款者更有可能按时偿还贷款。
	\item Condition: 包含了各方面的情况，比如贷款者所在区域的整体经济情况，贷款者所在行业的发展状况。
	\item Capacity: 判断贷款公司是否有足够的规模产生足够偿还贷款的现金流。
	\item Collateral: 抵押品，判断公司是否有足够价值的抵押品可供违约的时候进行偿还。
\end{itemize}
同时银行也可以通过客观的数据来判断贷款者的信用情况，比如通过具体的数据建立模型，判断贷款者的还款意愿。同时在贷款已经发放之后，银行还会追踪贷款者的各方面行为，包括：
\begin{itemize}
	\item 对贷款者信用状况的重新评估
	\item 对贷款使用的监控，防止贷款者将贷款用于不良的目的
	\item 对贷款者经营活动的监督
\end{itemize}
\subsection{基于统计学以及机器学习的风控方法}
贷款者一般都有很多可量化的客观指标，比如贷款者的年龄、性别、收入、家庭人口数量、贷款金额、贷款期限等。国外已经有大量研究采用不同模型对其进行拟合，从而预估贷款者的违约可能性。比如: 
\begin{itemize}
	\item 通过主成分分析，将贷款者通过聚类划为数个分组然后再分别进行线性拟合\cite{2016konocreditrisk}
	\item 用企业数十种的财务数据，通过神经网络判断其债务违约概率\cite{angelini2008neural}
\end{itemize}
\subsection{基于大数据的风控方法}
之前研究所使用的数据集往往规模较小，每个贷款者仅有几个到几十个特征。对于个人贷款，商业银行或互联网公司就很难搜集到足够的与贷款直接相关的财务信息。但是在大数据发达的今天，贷款提供方可以通过多种途径获得贷款者的其他信息，比如互联网金融公司的软件可以通过用户的手机权限，合法地获得用户的通话记录、地理位置等，从这些信息里可以提取大量用户的特征，以此建立模型。
\chapter{相关理论综述}
本文所要完成的任务是根据用户的数据来预测用户的违约概率，因此是一个回归任务。下面介绍了几种最广泛使用的回归分析的模型，并且分析了其评估的效果和可行性，从中选择了三种模型作为本文主要采用的预测模型。
\section{问题描述}
我们的目标是给定一组用户的特征$ x \in R^m $，求出该用户的违约概率 $ y \in [0, 1]$。
为了统一描述，对于某个风控的数据集，我们把第j个用户的特征记为$\vec x_j = [x_{j,1}, x_{j,2}, ..., x_{j,n}]$, 因此第j个用户的第i个特征就是$x_{j,i}$。第j个用户的标签为$y_j$, 满足
\[
y_j =\begin{cases}  0 \quad \text{用户j的贷款按时偿还} \\ 1 \quad \text{用户j的贷款违约} \end{cases}
\]

我们希望可以找到函数$f(\vec x)$ 使得对于每个用户j，$f(\vec x_j)$ 都尽量接近 $y_j$， 因此可以用均方误差作为损失函数，即优化目标为 $\min \limits_{f} \sum_j (f(\vec x_j) - y_j)^2$。可以使用如下所述的多种回归模型对其拟合。

\section{回归模型}

\textbf{线性回归}:
线性回归也称作最小二乘法拟合，是最简单的回归方法，最小化目标为
\[\min \limits_{w_i, b} \sum_{j} (\vec w \cdot \vec x_j + b - y_j)^2\]
 
通过最小二乘法可以对$\vec w$和$b$进行估计。线性回归易于实现，并且具有良好的可解释性，拟合出来的$w_i$就是第i个特征的权重。

\textbf{逻辑回归}:
逻辑回归是一种常用的回归方法。用均方误差作为损失函数的逻辑回归的目标为：
\[\min \limits_{w_i, b} \sum_{j} (\text{sigmoid} ( \vec w\cdot \vec x_j + b) - y_j)^2\]

其中，$ \text{sigmoid}(x) = \dfrac{1}{1+e^{-x}} $。逻辑回归一般使用梯度下降法求解最佳参数。

逻辑回归和线性回归一样易于实现，并具有良好的可解释性。所不同的是逻辑回归拟合的值域是[0,1]，更加适合二分类数据。

\textbf{支持向量机}:
支持向量机\cite{scholkopf2001svm}是一种强大的二分类算法，最小化目标为：
\[ \min \limits_{w_i, b} \dfrac{1}{n} \sum_{j}(\max(0, 1-y_j(\vec w \cdot x_j))) + \lambda ||w||^2  \]

其原理就是找到一个超平面能够将$y=0$和$y=1$的两类样本尽可能分离开。

支持向量机可以通过求解对偶问题把原问题转换为带约束的二次规划问题，然后通过多种方法求解。同时支持向量机也支持Kernel Tricks，即通过定义Kernel Function(核函数)把原数据映射到无穷维内积空间，然后进行拟合。常用的有高斯核等。支持向量机具有非常强大的二分类功能。

\textbf{人工神经网络}:
人工神经网络是一种从人类神经元连接方式启发而得到的机器学习模型，其每一层的基本方法是通过矩阵乘法将m维的数据映射到n维，然后再用激活函数产生输出，公式如下：
\[ \bm{y} = \text{activation}(W\bm{x} + \bm{b}) \]

神经网络可以有很多层，也可以有不同的结构，比如卷积神经网络、循环神经网络等。\cite{nndp} 神经网络具有强大的拟合能力，理论上足够大的单隐层神经网络可以逼近任意函数\cite{cybenko1989approximation}。神经网络在各种回归和分类任务上都具有良好的表现，但是容易出现过拟合问题，并且缺少可解释性。

\textbf{决策树}:
决策树\cite[311]{james2013introduction}是一种简单的模型，通过树结构来进行分类。决策树的每一个节点相当于一个if-else语句，对某个属性进行划分。对于该属性的值小于节点的值的样本，就把其送入左边的节点继续判断；反之则送入右边的节点。决策树把样本的向量空间分割为多个高维长方体，任何一个样本必然落在其中的一个长方体中。

决策树每一步从某一个叶节点出发，挑选一个属性进行分裂。一般来说分类决策树挑选分裂点的原则是信息增益，也就是让每个子节点的信息熵的和尽可能变小。也就是说其优化目标是：

\[ \min \sum_k \sum_m p_{mk}\log(p_{mk}) \]

其中， $p_{mk}$表示在叶子（区域）m上类别为k的样本的比例。

由于决策树的性质，理论上决策树可以对任何分布的样本进行分类。但是也很容易出现过拟合的现象，所以现在广泛采用大量浅层决策树的集成模型进行预测，如下面介绍的随机森林和XGBoost。而且决策树具有非常好的可解释性，因此被广泛应用于各种特征无关的分类工作。

\textbf{随机森林}:
随机森林\cite[320]{james2013introduction}是多个决策树的集成模型，通过Boostrap采样算法从原始数据集中生成多个新的数据集，并分别用其训练不同的决策树模型。其中决策树选择分裂的时候仅仅从所有属性里面的一个子集挑选属性，以此降低过拟合的问题。
随机森林能够显著提高决策树的效果，并且保留了决策树的可解释性。

\textbf{XGBoost}:
XGBoost\cite{chen2016xgboost}是另一种决策树的集成模型，可以用于回归和分类。对于某一样本，其输出就是XGBoost模型中的每一棵决策树对应的叶子节点的权重之和。XGBoost通过Boosting算法，在每轮迭代中加入一棵新的决策树来尽可能降低整体的损失函数。XGBoost通过限制单棵决策树的深度、部分采样以及对叶子和权重的规范化来减少过拟合。

XGBoost的算法通过预先索引、对属性的并行，部分采样等方式提高每棵树的生成速度，极大提高了运算效率。
\\

因为我们的问题是对用户是否违约进行预测，是一个二分类问题，因此直接使用线性回归的效果较差。支持向量机的运算开销过大，需要十几个小时才能做出一组预测，因此在之后的实验中没有采用该方法。同时决策树、随机森林和Xgboost呈现递进关系，Xgboost的一轮迭代就是决策树。因此在后文中我们采取了其中逻辑回归、人工神经网络和Xgboost这三个模型进行实验。

\section{特征处理方法}
如果原始数据的维度过高，或是存在一些缺失值，会影响模型的拟合效果，或是增加拟合的成本。一般可以通过主成分分析解决原始数据维度过高的问题，通过AutoEncoder去除原始数据中的缺失值、噪声等异常数据。

\textbf{主成分分析}:
主成分分析\cite{shlens2014pca} 是一种降维的方法，把样本空间从m维降到k维(k<m)，满足
\[
\min \sum_{i = 1}^{n} ||\bm{x_i} - \sum_{j=1}^{k}(<\bm{x_i},\bm{\alpha_j}>\alpha_j)||^2
\]

即当样本从k维的“压缩状态”回到m维的时候，其损失(用差值的模平方衡量)之和最小。因此这是一种保留最多信息的线性映射。

\textbf{AutoEncoder}:
AutoEncoder是一种输入和输出的大小相同的神经网络，其原理是把N维的输入数据转换到m维的隐藏层(m<n)，然后再转换成N维。AutoEncoder可以通过提取隐藏层来实现数据降维，也可以通过再设置一个与输入同样大小的输出层，对输入数据进行数据缺失值的填补。\cite{vincent2008autoencoder}

因为现在深度学习的任务都可以在GPU上运行，因此可以接受维度很高的数据，而主成分分析必然损失一部分信息。因此在之后的讨论中我们均使用原始数据，并未采用主成分分析对其进行降维。

\section{模型效果评估方法}
\subsection{训练集和测试集}
为了评估模型的预测效果，我们需要从原始数据中划分一部分数据集作为测试集，测试集不参与训练。因为各种深度学习模型都具有强大的拟合能力，所以模型在训练集上的误差会趋向于0，所以为了评估模型真实的拟合能力，我们需要让训练好的模型在测试集上做出预测，查看其效果。
\subsection{K折交叉验证}
随着现代计算机计算性能的提升，对于大数据集的计算也不再昂贵。因此现在一般采用K折交叉验证的方式评估模型性能。即，把数据集分成K等份，依次选择其中的K-1份作为训练集，剩下的一份作为验证集，分别计算其误差。通过这种方式，可以显著降低评估误差，降低某些偶然因素导致的评估的不准确性，找出最佳的模型。\cite[213]{james2013introduction}
\subsection{模型评估指标}
\subsubsection{准确率}
模型的输出结果如果大于等于0.5，就认为其预测的标签是1，反之则为0。则准确率计算公式如下: 
\[
Accuracy = \dfrac{\text{预测准确的样本数}}{\text{总样本数}}
\]
\subsubsection{AUC(Area Under Roc Curve)}
ROC曲线，根据预测样本的真阳性率和假阳性率形成的曲线。其计算公式为
\[
    AUC = \int_{fpr=0}^1 tpr \mathrm{d}tpr
\]
其中，fpr表示假阳性率（实际标签为0，预测标签为1的样本数量/实际标签为0的样本总数（正样本总数））,tpr表示真阳性率（实际标签为1，预测标签为1的样本数量/实际标签为1的样本总数（负样本总数））
我们令预测标签为1的概率门槛从1变到0，则真阳性率和假阳性率将会从0变到1。设当有i个假阳性样本时，则假阳性率为i/N，N为所有实际标签为0的样本。不妨设第i的假阳性样本是刚刚加入的，也就是说模型对其的输出是这i个里面最小的，计其为$p_i$，则此时对应的真样本的数目为$\sum\limits_{j \in \text{正样本}} 1_{f(j) > f(i)}$
其中，
\[
1_{x>y} =\begin{cases}  0 \quad x<y \\ 1 \quad x>y \end{cases}
\]
根据以上的推导，可以看出AUC可以用如下公式计算：
\[
    AUC = \sum\limits_{i \in \text{正样本} } \sum\limits_{j \in \text{负样本} }{1_{f(i)>f(j)}}
\]
随机预测的AUC的值在0.5左右，AUC的值越大，说明模型越倾向于使得正样本的输出大于负样本的输出。
\chapter{风控原始数据预处理}
本文中所使用的数据集来自某互联网公司的数个风控平台的数据库。因为原始数据的格式不统一，含有大量字符串、压缩数据、列表、JSON等数据，因此必须经过处理才可以作为模型输入。数据预处理的好坏会极大影响模型预测的效果\cite{kotsiantis2006data}。我们的处理目标就是把原始数据中可能有用的数据尽可能提取成数值，同时尽可能丢弃无效的数据，保留对模型预测有效的数据，返回一个 (用户数目$\times$单个用户特征数目大小) 的表。并且表中的每一项数值都介于0到1之间，有利于对于数值敏感的模型的训练。
\section{大数据平台Hadoop}
Hadoop\footnote{https://hadoop.apache.org}是一个开源的大数据平台，提供了一套分布式存储和分布式任务调度的框架，能够让服务器集群更加有效地处理超大规模的数据。Spark\footnote{https://spark.apache.org/docs/latest/}是一套快速通用的集群计算系统，能够执行大量数据处理的任务，并且提供了Python、Java和Scala语言的高级API。Spark可以通过Hadoop进行任务调度。
\section{数据格式统一}
\subsection{JSON数据的提取}
JSON\footnote{https://www.json.org}是一种轻量级的数据表示格式，可以表示key-value类型的对象，以及列表，支持嵌套。每个对象的值可以是对象、字符串、数字或者是布尔值以及null。下列代码是一个简单的JSON数据示例。
\begin{lstlisting}[language=PYTHON, caption={JSON示例}, label={JSON示例}]
{
	"姓名":"张三",
	"家属": 
	{
		"父亲":"张父",
		"母亲":"张母"
	},
	"出生日期":"1999-9-9",
	"学历记录":
	[
		{
			"中学":"张村中学",
			"成绩":"95"
		},
		{
			"大学":"中国科学技术大学",
			"成绩":"1.0"
		}
	]
}
\end{lstlisting}

对于JSON格式的数据，处理方式为:

递归地取出所有非列表结构的字段。为了在数据处理之后不丢失原有的属性的信息，按照属性名拼接的方式产生列名。比如对于上面的示例中的“父亲”一项，就把值保存在列名为"家属\_\_\_父亲"的列中。对于列表形式的数据，根据实际情况，一般列表当中的对象都是字典类型，所以把其中所有数值类型字段进行求和和平均，加入列明为属性名拼接加上“\_\_\_SUM”或"\_\_\_AVG"标识符的列中，比如上面示例里的“学历记录”可以产生两列名为“学历记录\_\_\_成绩\_\_\_SUM”和“学历记录\_\_\_成绩\_\_\_AVG”的记录。

\begin{algorithm}[htbp]
	\SetAlgoLined
	\KwData{Json}
	\KwResult{OutputDict}
	OutputDict = \{\}\\
	UnParsedJsonElems = [Json]\\
	Prefix = ""\\
	\While{UnParsedJsonElems is not empty}{
		elem = UnParsedJsonElems.pop(0)\\
		\For{attr in elem.stringAndNumericalAttributes}
		{
			OutputDict[Prefix + attr] = elem[attr]\\
		}
		\For{attr in elem.listAttributes}
		{
			\For{subAttr in elem[attr][0].numericalAttributes}
			{
				OutputDict[Prefix+"\_\_\_"+attr + "\_\_\" + subAttr + "\_\_\_SUM"] = SUM(elem[attr], subAttr) \\
				OutputDict[Prefix+"\_\_\_"+attr + "\_\_\" + subAttr + "\_\_\_AVG"] = AVG(elem[attr], subAttr) \\
			}
		}
		\For{attr in elem.objectAttributes}
		{
			UnparsedJsonElem.push(elem[attr])
		}
	}
	\caption{Json处理算法示例}
	\label{algo:algorithm1}
\end{algorithm}
\subsection{特大数量异构JSON数据融合}
由于实验的数据集中含有数十万条用户的JSON数据，而且通过样本得知每个样本的JSON数据都略有不同，比如用户2可能有用户1中不存在的字段，用户1也可能有用户2中不存在的字段。由于整体数据量过于庞大，如果要遍历所有的用户数据，由于各用户JSON字段的差别，无法并行地进行处理，因此效率极其低下。为了获取大部分用户都拥有的字段，我们采用采样的方法。

假设某一字段a, 在数据集的所有用户中出现的频率是p，那么缺失率就是1-p。在N次采样中，字段a都没有出现的概率为
\[ p_{\text{miss in N samples}} = (1-p)^N\]

因此如果取采样数目为80，那么对于在样本中出现频率仅为10\%的字段，其在采样中不出现的概率也仅仅不到$\frac{1}{5000}$，因此我们可以认为大小为80的采样样本中包含了几乎所有出现频率大于0.1的字段。

因此首先把采样中出现的字段存放在字典中，将每个字段对应到一个整数，可以大大减少因为字段过长产生的冗余内存。然后通过该字典把用户的JSON数据进行统一处理，如果用户含有不包含在字典中的JSON字段则丢弃，如果用户缺少字典中的JSON字段，则将该字段补为空值。同时把字段名都改为字典中对应字段的整数Value。

\subsection{多源数据整合}
数据集中的每个表都有一个唯一的用户ID标识，因此可以通过表连接的方式获得总表。
\subsection{数据的数值化转换}
由于不同的数据源使用不同的方法表示数据的缺失、异常值，比如有的采用字符串"无数据"表示，有的使用null值表示，有的使用空字符串表示，所以转换的时候需要注意，并且将其都转换为nan值。

为了获取每一列可能的非数值数据，通过遍历的方法，逐列进行检查，把每一列出现的非数值值记录到文件里并且人工检查之。尽可能把能够量化的字符串变量进行数值化转换。如对于“用户与律师通话频率”字段，有“频繁通话”、“正常通话”、“很少通话”、“从未通话”四种字符串，可以分别转换为1、0.667、0.333、0四种数值。

对于类别类型的数据，比如“省份”之类，采用one-hot编码，也就是说把列数从原来的一列扩充到与类别数量相同的列数。



\section{数据清洗}
由于通过前面步骤得到的数据量依然非常巨大，而大量对模型预测无效的数据会对模型产生不良影响，同时数据量过大也会极大地增加模型的训练和运行成本，所以我们需要对每一列数据进行初步分析，删除那些可能无效的列。

\subsection{数据的初步分析}
我们采用的方法是对每一列数据进行分析。
\begin{itemize}
	\item 缺失率：即该列数据在所有样本中缺失的比率。如果缺失率过高，则考虑丢弃该列。
	\item 最大值、最小值、标准差。这些数值可以反映该列数据的变动程度。如果最大值非常接近最小值或者标准差过小，则可认为该列数据在每一个样本上变动很小，没有反映出样本之间的差异性，考虑丢弃该列。
	\item 与标签的相关系数。如果与标签的相关系数很小，则说明该列很可能没有有用信息。但是也有可能存在该列数据仅仅和标签的一阶相关系数小，可能有其他的非线性关系或者和别的特征组合影响标签的情况。
\end{itemize}
我们通过以上三点综合考虑每一列数据的有效性，最终决定是否要丢弃该列数据。

为了使得最后的数据分布在0—1之间，把每一列都按照如下方式处理：

$normalized\_value = \dfrac{value - min}{max - min}$
\input{chapters/my_data_imputation.tex}
\chapter{模型测试与评估}
本文需要进行三组实验，分别是
\begin{enumerate}
    \item 对是否违约的标签的预测：用第\ref{data_refinement}章处理之后的数据，将其跟之前的研究\cite{czh}处理的数据，使用三种模型进行预测，对比显示出通过更加细致的数据处理工作产生的数据对模型预测效果的提高。
    \item 对缺失值填补算法的效果评估和对比：使用第\ref{data_refinement}章处理后的数据，将第\ref{data_imputation}章提出的基于嵌入的缺失值填补方法和广泛采用的AutoEncoder进行对比。
    \item 缺失值填补对预测效果的影响：把缺失值填补之后的数据和填补之前的数据，用三种模型分别进行预测，判断缺失值填补对于模型预测效果的影响。
\end{enumerate}
\section{对是否违约标签的预测}
\subsection{数据设置}
我们从前期整理的风控数据集里取出100000条用户信息，然后使用5折交叉验证方法，每次用80000条用户信息作为训练集，用20000条用户信息作为测试集。

对预测数据的评估方法采用准确率和AUC分数。准确率用0.5作为分界线。

对于无法处理缺失值的模型，比如神经网络，采用缺失值补0的方法进行预测。
\subsection{模型设置}
由于线性模型的拟合能力显然低于逻辑回归，而支持向量机对此类型大数据的训练相当缓慢。而随机森林的效率和准确度被认为是不如其改进版XGBoost的。因此在这里我们仅采用逻辑回归、神经网络以及XGBoost三种模型进行测试。

对于逻辑回归，我们把其当作单层神经网络，和神经网络一起使用Tensorflow框架\footnote{tensorflow.org}进行训练。XGBoost模型则使用python的xgboost库\footnote{xgboost.org}。
\begin{itemize}
    \item 逻辑回归：学习率0.1，使用梯度下降优化器。损失函数为均方误差。
    \item 神经网络：隐层神经元数目为100，学习率0.002，使用Adam优化器，激活函数全部使用sigmoid。损失函数为均方误差。
    \item XGBoost：树深度为5，学习率(eta)0.3，目标：二分类逻辑回归，列采样率0.6。
\end{itemize}
模型中若有为声明的其他参数，均按照库函数的缺省值为准。我们对所有模型都进行5折交叉验证，迭代40次。然后分别以AUC Score和准确率进行评判。
\subsection{实验结果}
\begin{table}[htp]
    \centering\small
    \caption{模型效果对照表}
    \begin{tabular}{lllllllll}
    \toprule
     & AUC-Score              &              &              &  & Accuracy              &              &             &  \\
    \midrule
    epoch & LR              & DNN             & XGB             &  & LR              & DNN             & XGB             &  \\
    ... & ...             & ...             & ...             &     & ...             & ...             & ...             &  \\
    10  & 0.8147          & 0.8419          & 0.8599          &     & 0.7782          & 0.7949          & 0.7963          &  \\
    11  & 0.8152          & 0.8429          & 0.8614          &     & 0.7794          & 0.7942          & 0.7978          &  \\
    12  & 0.8154          & 0.8434          & 0.8626          &     & 0.7786          & 0.7924          & 0.7993          &  \\
    13  & 0.8153          & 0.8442          & 0.8637          &     & 0.7794          & 0.7934          & 0.8001          &  \\
    14  & 0.8158          & \textbf{0.8459} & 0.8648          &     & 0.7797          & \textbf{0.7981} & 0.8009          &  \\
    15  & 0.8160          & 0.8444          & 0.8657          &     & 0.7792          & 0.7941          & 0.8015          &  \\
    ... & ...             & ...             & ...             &     & ...             & ...             & ...             &  \\
    30  & 0.8177          & 0.8315          & 0.8767          &     & 0.7802          & 0.7867          & 0.8111          &  \\
    31  & 0.8175          & 0.8338          & 0.8771          &     & \textbf{0.7815} & 0.7880          & 0.8113          &  \\
    32  & 0.8178          & 0.8309          & 0.8776          &     & 0.7809          & 0.7867          & 0.8120          &  \\
    ... & ...             & ...             & ...             &     & ...             & ...             & ...             &  \\   
    37  & 0.8174          & 0.8280          & 0.8797          &     & 0.7807          & 0.7854          & 0.8147          &  \\
    38  & \textbf{0.8185} & 0.8275          & 0.8800          &     & 0.7782          & 0.7855          & 0.8155          &  \\
    39  & 0.8183          & 0.8238          & \textbf{0.8805} &     & 0.7804          & 0.7807          & \textbf{0.8157} &  \\
    \bottomrule
    \end{tabular}
\end{table}
可见在LR、DNN、Xgboost三种模型中，Xgboost的表现要明显优于前两者。通过与同一数据集之前的测试结果\cite{czh}进行比较，可以看出，通过第 \ref{data_refinement} 章中采用的数据处理方法得到的数据，在使用同样的模型预测时，效果要显著好于之前的结果。

\begin{table}[htb]
\centering\small
\caption{前后结果对照}
    \begin{tabular}{ccc}
        \toprule
        模型-指标  & 之前的数据 & 本文处理的数据 \\
        \midrule
        LR-AUC  & 0.770  & 0.8185 \\
        DNN-AUC & 0.779  & 0.8469 \\
        XGB-AUC & 0.785  & 0.8805 \\
        LR-ACC & 0.711  & 0.7815 \\
        DNN-ACC & 0.717  & 0.7981 \\
        XGB-ACC & 0.721  & 0.8157 \\
        \bottomrule
    \end{tabular}
\end{table}

实验结果证明了对数据的预处理在很大程度上能够影响模型的预测效果。之前的数据处理基本上只是提取出了数值数据，没有考察很多字符串类型的类别数据，同时把很多“非数值”类型的“伪数值”数据也当作数值传入（比如手机号码的各种数据），增大了样本的噪声。而本文实行的数据处理人工考察了数据，并且对列表类型的数据丢弃了最大最小值，加入了列表总长。


\section{对缺失值填补算法的效果评估和对比}
\subsection{数据设置}
我们使用第\ref{data_refinement}章中处理好的风控数据来产生新的三元组数据集。

因为原数据集含有100,000个样本，每个样本有1156个特征，因此总共可以转化成115,600,000条三元组数据。其中缺失的数据总共有17,045,944个，缺失率为14.7\%。

因为三元组数据的数据量非常巨大，因此只需要用一小部分作为测试集就可以测出较为准确的效果。我们把所有三元组数据中的0.03\%作为测试集，测试集含有30000多条三元组数据。

\subsection{模型设置}
\begin{itemize}
    \item AutoEncoder:  使用单层的AutoEncoder，神经元数量为256，激活函数为sigmoid，训练Batch大小为100。
    \item 特征嵌入模型：样本和特征的嵌入向量维度都是128。对于嵌入向量加入规范化项，权重为0.1。内部神经网络的隐层神经元数为64，隐层激活函数为relu，输出层激活函数为sigmoid。
\end{itemize}

所有模型均采用学习率为0.02的Adam Optimzer执行梯度下降算法。损失函数使用均方误差。
\subsection{实验结果}
\begin{table}[htb]
    \caption{模型效果对照表}
    \centering
    \begin{tabular}{lll}
    \toprule
    Epochs & AutoEncoder & 特征嵌入模型 \\
    \midrule
    1 & 0.00845 & 0.00479 \\
    2 & \textbf{0.00834} & 0.00401 \\
    ... & ... & ... \\
    6 & 0.01451 & 0.00327 \\
    7 & 0.01555 & \textbf{0.00319} \\
    8 & 0.01680 & 0.00336 \\
    \bottomrule
    \end{tabular}
\end{table}
相比于AutoEncoder，使用特征嵌入模型的缺失值填补要更为精准，在数次迭代之后，AutoEncoder出现了过拟合现象，但是特征嵌入模型保持着较低的损失。

\section{缺失值填补对预测效果的影响}
由于我们使用的原始数据集是多个数据集进行内联得到的，所以数据的稀疏性并不是很大，缺失比率不到15\%，所以在多次实验中发现使用缺失值填补算法并不能对预测效果有改善，甚至因为引入了新的噪声会使预测效果出现微小的下降。

但是如果把数据集中的一半非缺失数据丢弃，使得数据缺失比率达到60\%，则LR回归和神经网络的预测效果会因为缺失值的填补而得到提升，但是Xgboost的预测效果依然出现了下降。
\subsection{数据设置}
实验数据集分为未填补缺失值的原始数据集和经过缺失值填补的数据集。对于每一数据集都采用5折交叉验证，最多进行40轮迭代，并且从其中选出平均分数最高的一轮，进行对比。
\subsection{模型设置}
缺失值填补模型采用上述的特征嵌入模型，参数与上述实验中的模型参数相同，回归模型则采用第一节中各个用于原始数据回归的LR、DNN、Xgb模型，参数亦保持不变。
\subsection{实验结果}
\begin{table}[htb]
    \caption{模型效果对照表}
    \centering
    \begin{tabular}{lll}
    \toprule
    模型-指标 & 缺失值未填补 & 缺失值填补 \\
    \midrule
    LR-auc & 0.7454 & 0.7591 \\
    LR-acc & 0.7545 & 0.7580 \\
    DNN-auc & 0.7590 & 0.7648 \\
    DNN-acc & 0.7590 & 0.7600 \\
    Xgb-auc & 0.8275 & 0.8201 \\
    Xgb-acc & 0.7776 & 0.7730 \\
    \bottomrule
    \end{tabular}
\end{table}

可见Xgboost在缺失值未填补的情况下效果最好。我们猜测这是因为Xgboost自身可以动态选择把缺失值划归哪一颗子树，因此能够较好地处理缺失值。而神经网络等算法自身并不能处理缺失数据，必须用一个值进行代替（比如是0或平均值，本实验中用的是0），因此缺失率过高时，会学习到“错误信息”，因此填补缺失值能够提高预测效果。
\input{chapters/my_conclusion_and_expectation.tex}
%% !TeX root = ../main.tex

\chapter{简介}

\section{一级节标题}

\subsection{二级节标题}

\subsubsection{三级节标题}

\paragraph{四级节标题}

\subparagraph{五级节标题}

Lorem ipsum dolor sit amet, consectetur adipiscing elit, sed do eiusmod tempor
incididunt ut labore et dolore magna aliqua.
Ut enim ad minim veniam, quis nostrud exercitation ullamco laboris nisi ut
aliquip ex ea commodo consequat.
Duis aute irure dolor in reprehenderit in voluptate velit esse cillum dolore eu
fugiat nulla pariatur.
Excepteur sint occaecat cupidatat non proident, sunt in culpa qui officia
deserunt mollit anim id est laborum.



\section{脚注}

Lorem ipsum dolor sit amet, consectetur adipiscing elit, sed do eiusmod tempor
incididunt ut labore et dolore magna aliqua.
\footnote{Ut enim ad minim veniam, quis nostrud exercitation ullamco laboris
  nisi ut aliquip ex ea commodo consequat.
  Duis aute irure dolor in reprehenderit in voluptate velit esse cillum dolore
  eu fugiat nulla pariatur.}

%% !TeX root = ../main.tex

\chapter{浮动体}

\section{三线表}

三线表是《撰写手册》推荐使用的格式，如表~\ref{tab:exampletable}。
\begin{table}[htb]
  \centering\small
  \caption{表号和表题在表的正上方}
  \label{tab:exampletable}
  \begin{tabular}{cl}
    \toprule
    类型   & 描述                                       \\
    \midrule
    挂线表 & 挂线表也称系统表、组织表，用于表现系统结构 \\
    无线表 & 无线表一般用于设备配置单、技术参数列表等   \\
    卡线表 & 卡线表有完全表，不完全表和三线表三种       \\
    \bottomrule
  \end{tabular}
  \note{注：表注分两种，第一种是对全表的注释，用不加阿拉伯数字排在表的下边，
    前面加“注：”；第二种是和表内的某处文字或数字相呼应的注，
    在表里面用带圈的阿拉伯数字在右上角标出，然后在表下面用同样的圈码注出来}
\end{table}

编制表格应简单明了，表达一致，明晰易懂，表文呼应、内容一致。
排版时表格字号略小，或变换字体，尽量不分页，尽量不跨节。
表格太大需要转页是，需要在续表上方注明“续表”，表头页应重复排出。



\section{插图}

有的同学可能听说“\LaTeX{} 只能使用 eps 格式的图片”，甚至把 jpg 格式转为 eps。
事实上，这种做法已经过时。
而且每次编译时都要要调用外部工具解析 eps，导致降低编译速度。
所以我们推荐矢量图直接使用 pdf 格式，位图使用 jpeg 或 png 格式。
\begin{figure}[htb]
  \centering
  \includegraphics[width=0.3\textwidth]{ustc_logo_fig.pdf}
  \caption{图号、图题置于图的下方}
  \label{fig:logo}
  \note{注：图注的内容不宜放到图题中。}
\end{figure}

关于图片的并排，推荐使用较新的 \pkg{subcaption} 宏包，
不建议使用 \pkg{subfigure} 或 \pkg{subfig} 等宏包。



\section{算法环境}

模板中使用 \pkg{algorithm2e} 宏包实现算法环境。关于该宏包的具体用法，
请阅读宏包的官方文档。

\begin{algorithm}[htb]
  \small
  \SetAlgoLined
  \KwData{this text}
  \KwResult{how to write algorithm with \LaTeX2e }

  initialization\;
  \While{not at end of this document}{
    read current\;
    \eIf{understand}{
      go to next section\;
      current section becomes this one\;
    }{
      go back to the beginning of current section\;
    }
  }
  \caption{算法示例1}
  \label{algo:algorithm1}
\end{algorithm}

注意，我们可以在论文中插入算法，但是插入大段的代码是愚蠢的。
然而这并不妨碍有的同学选择这么做，对于这些同学，建议用 \pkg{listings} 宏包。

%% !TeX root = ../main.tex

\chapter{数学}

\section{数字和单位}

宏包 \pkg{siunitx} 提供了更好的数字和单位支持：
\begin{itemize}
  \item \num{12345.67890}
  \item \num{1+-2i}
  \item \num{.3e45}
  \item \num{1.654 x 2.34 x 3.430}
  \item \si{kg.m.s^{-1}}
  \item \si{\micro\meter} $\si{\micro\meter}$
  \item \si{\ohm} $\si{\ohm}$
  \item \numlist{10;20}
  \item \numlist{10;20;30}
  \item \SIlist{0.13;0.67;0.80}{\milli\metre}
  \item \numrange{10}{20}
  \item \SIrange{10}{20}{\degreeCelsius}
\end{itemize}



\section{数学符号和公式}

\LaTeX{} 默认按照美国的习惯排版数学公式和符号，
但是《撰写手册》要求数学符号依据《GB 3102.11--1993》执行，
与 \LaTeX{} 的习惯有所差异。
本模板基于 \pkg{unicode-math} 配置数学符号，以遵循国标的规定。

注意，\pkg{unicode-math} 宏包与 \pkg{amsfonts}, \pkg{amssymb}, \pkg{bm},
\pkg{mathrsfs}, \pkg{upgreek} 等宏包\emph{不}兼容。
本模板作了处理，用户可以直接使用这些宏包的命令，如 \cs{bm}, \cs{mathscr},
\cs{upGamma}。

本模板中数学符号的用法与 \LaTeX{} 传统有些区别：
\begin{itemize}
  \item 数学常数和特殊函数使用正体，
    如圆周率 $\symup{\pi}$、$\symup{\Gamma}$ 函数。
    应使用 \pkg{unicode-math} 宏包提供的 \cs{symup} 命令转为正体，
    如 \verb|\symup{\pi}|。
  \item 向量和矩阵粗斜体，应使用 \cs{symbf} 命令，
    如 \verb|\symbf{u}|、\verb|\symbf{A}|。
  \item 有限增量符号 $\increment$ （U+2206）应使用 \cs{increment} 命令。
  \item 微分符号 $\dif$ 使用正体，本模板提供了 \cs{dif} 命令。
\end{itemize}

除此之外，模板还提供了一些命令方便使用：
\begin{itemize}
  \item 常数 $\upe$：\verb|\upe|
  \item 负数单位 $\upi$：\verb|\upi|
  \item 圆周率 $\uppi$：\verb|\uppi|
  \item $\argmax$：\verb|\argmax|
  \item $\argmin$：\verb|\argmin|
\end{itemize}

关于数学符号更多的用法，参见 \pkg{unicode-math} 宏包的使用说明和符号列表
\pkg{unimath-symbols}。

在编辑数学公式时，最好避免直接使用字体命令，
而应该定义一些语义命令取代字体命令，
这样输入更简单，也让 \LaTeX{} 代码更有可读性，
而且还方便根据需要统一修改改格式，比如：
\begin{itemize}
  \item 向量 $\vec{x}$：\verb|\renewcommand\vec{\symbf}|
  \item 矩阵 $\mat{A}$：\verb|\newcommand\mat{\symbf}|
  \item 张量 $\ts{T}$： \verb|\newcommand\ts{\symbfsf}|
\end{itemize}

更多的例子：
\begin{equation}
  \upe^{\upi\uppi} + 1 = 0
\end{equation}
\begin{equation}
  \frac{\dif^2 u}{\dif t^2} = \int f(x) \dif x
\end{equation}
\begin{equation}
  \argmin_x f(x)
\end{equation}
\begin{equation}
  \mat{A} \vec{x} = \lambda \vec{x}
\end{equation}



\section{定理和证明}

示例文件中使用 \pkg{amsthm} 宏包配置了定理、引理和证明等环境。
用户也可以使用  \pkg{ntheorem} 宏包。

\begin{definition}
  If the integral of function $f$ is measurable and non-negative, we define
  its (extended) \textbf{Lebesgue integral} by
  \begin{equation}
    \int f = \sup_g \int g,
  \end{equation}
  where the supremum is taken over all measurable functions $g$ such that
  $0 \le g \le f$, and where $g$ is bounded and supported on a set of
  finite measure.
\end{definition}

\begin{assumption}
The communication graph is strongly connected.
\end{assumption}

\begin{example}
  Simple examples of functions on $\real^d$ that are integrable
  (or non-integrable) are given by
  \begin{equation}
    f_a(x) =
    \begin{cases}
      |x|^{-a} & \text{if } |x| \le 1, \\
      0        & \text{if } x > 1.
    \end{cases}
  \end{equation}
  \begin{equation}
    F_a(x) = \frac{1}{1 + |x|^a}, \qquad \text{all } x \in \real^d.
  \end{equation}
  Then $f_a$ is integrable exactly when $a < d$, while $F_a$ is integrable
  exactly when $a > d$.
\end{example}

\begin{lemma}[Fatou]
  Suppose $\{f_n\}$ is a sequence of measurable functions with $f_n \geq 0$.
  If $\lim_{n \to \infty} f_n(x) = f(x)$ for a.e. $x$, then
  \begin{equation}
    \int f \le \liminf_{n \to \infty} \int f_n.
  \end{equation}
\end{lemma}

\begin{remark}
  We do not exclude the cases $\int f = \infty$,
  or $\liminf_{n \to \infty} f_n = \infty$.
\end{remark}

\begin{corollary}
  Suppose $f$ is a non-negative measurable function, and $\{f_n\}$ a sequence
  of non-negative measurable functions with
  $f_n(x) \le f(x)$ and $f_n(x) \to f(x)$ for almost every $x$. Then
  \begin{equation}
    \lim_{n \to \infty} \int f_n = \int f.
  \end{equation}
\end{corollary}

\begin{proposition}
  Suppose $f$ is integrable on $\real^d$. Then for every $\epsilon > 0$:
  \begin{enumerate}
    \renewcommand{\theenumi}{\roman{enumi}}
    \item There exists a set of finite measure $B$ (a ball, for example) such
      that
      \begin{equation}
        \int_{B^c} |f| < \epsilon.
      \end{equation}
    \item There is a $\delta > 0$ such that
      \begin{equation}
        \int_E |f| < \epsilon \qquad \text{whenever } m(E) < \delta.
      \end{equation}
  \end{enumerate}
\end{proposition}

\begin{theorem}
  Suppose $\{f_n\}$ is a sequence of measurable functions such that
  $f_n(x) \to f(x)$ a.e. $x$, as $n$ tends to infinity.
  If $|f_n(x)| \le g(x)$, where $g$ is integrable, then
  \begin{equation}
    \int |f_n - f| \to 0 \qquad \text{as } n \to \infty,
  \end{equation}
  and consequently
  \begin{equation}
    \int f_n \to \int f \qquad \text{as } n \to \infty.
  \end{equation}
\end{theorem}

\begin{proof}
  Trivial.
\end{proof}

\newtheorem*{axiomofchoice}{Axiom of choice}
\begin{axiomofchoice}
  Suppose $E$ is a set and ${E_\alpha}$ is a collection of
  non-empty subsets of $E$. Then there is a function $\alpha
  \mapsto x_\alpha$ (a ``choice function'') such that
  \begin{equation}
    x_\alpha \in E_\alpha,\qquad \text{for all }\alpha.
  \end{equation}
\end{axiomofchoice}

\newtheorem{observation}{Observation}
\begin{observation}
  Suppose a partially ordered set $P$ has the property
  that every chain has an upper bound in $P$. Then the
  set $P$ contains at least one maximal element.
\end{observation}
\begin{proof}[A concise proof]
  Obvious.
\end{proof}

%% !TeX root = ../main.tex

\chapter{引用文献的标注}

模板使用 \pkg{natbib} 宏包来设置参考文献引用的格式，
更多引用方法可以参考该宏包的使用说明。



\section{顺序编码制}

\subsection{角标数字标注法}

\citestyle{super}
\noindent
\begin{tabular}{l@{\quad$\Rightarrow$\quad}l}
  \verb|\cite{knuth86a}|         & \cite{knuth86a}         \\
  \verb|\citet{knuth86a}|        & \citet{knuth86a}        \\
  \verb|\cite[42]{knuth86a}|     & \cite[42]{knuth86a}     \\
  \verb|\cite{knuth86a,tlc2}|    & \cite{knuth86a,tlc2}    \\
  \verb|\cite{knuth86a,knuth84}| & \cite{knuth86a,knuth84} \\
\end{tabular}


\subsection{数字标注法}

\citestyle{numbers}
\noindent
\begin{tabular}{l@{\quad$\Rightarrow$\quad}l}
  \verb|\cite{knuth86a}|         & \cite{knuth86a}         \\
  \verb|\citet{knuth86a}|        & \citet{knuth86a}        \\
  \verb|\cite[42]{knuth86a}|     & \cite[42]{knuth86a}     \\
  \verb|\cite{knuth86a,tlc2}|    & \cite{knuth86a,tlc2}    \\
  \verb|\cite{knuth86a,knuth84}| & \cite{knuth86a,knuth84} \\
\end{tabular}



\section{著者-出版年制标注法}

\citestyle{authoryear}
\noindent
\begin{tabular}{l@{\quad$\Rightarrow$\quad}l}
  \verb|\cite{knuth86a}|         & \cite{knuth86a}         \\
  \verb|\citep{knuth86a}|        & \citep{knuth86a}        \\
  \verb|\cite[42]{knuth86a}|     & \cite[42]{knuth86a}     \\
  \verb|\cite{knuth86a,tlc2}|    & \cite{knuth86a,tlc2}    \\
  \verb|\cite{knuth86a,knuth84}| & \cite{knuth86a,knuth84} \\
\end{tabular}

\vskip 2ex \citestyle{super}
注意，参考文献列表中的每条文献在正文中都要被引用
\cite{slg,lyc,ljs,cgw,cjb,kqy,yhs,yx,dwx,jxz,wjk,syw,wf,xd,twh,huston}。

\bibliography{bib/ustc}
\appendix
%% !TeX root = ../main.tex

\chapter{补充材料}


补充内容。


\backmatter

%% !TeX root = ../main.tex

\begin{acknowledgements}

在研究学习期间，我有幸得到了三位老师的教导，
他们是：我的导师，中国科大XXX研究员，中科院X昆明动物所马老师以及美国犹他大学的XXX老师。
三位深厚的学术功底，严谨的工作态度和敏锐的科学洞察力使我受益良多。
衷心感谢他们多年来给予我的悉心教导和热情帮助。

感谢XXX老师在实验方面的指导以及教授的帮助。
科大的XXX同学和XXX同学参与了部分试验工作，在此深表谢意。

\end{acknowledgements}

%% !TeX root = ../main.tex

\begin{publications}

\section*{已发表论文}

\begin{enumerate}
\item A A A A A A A A A
\item A A A A A A A A A
\item A A A A A A A A A
\end{enumerate}

\section*{待发表论文}

\begin{enumerate}
\item A A A A A A A A A
\item A A A A A A A A A
\item A A A A A A A A A
\end{enumerate}

\section*{研究报告}
\begin{enumerate}
\item A A A A A A A A A
\item A A A A A A A A A
\item A A A A A A A A A
\end{enumerate}

\end{publications}


\end{document}
